\documentclass[12pt]{article}
\pagestyle{empty}
\usepackage{amsmath, amssymb, amsthm}
\usepackage{latexsym, epsfig, ulem, cancel, multicol, hyperref}
\usepackage{graphicx, tikz, subfigure,pgfplots}
\usepackage[margin=1in]{geometry}
\setlength{\parindent}{0pt}
\usepackage{multirow}
\usepackage{mathtools}


\newcommand{\R}{\mathbb{R}}
\newcommand{\dydx}{\frac{dy}{dx}}
\usepackage{verbatim}
\usepackage{tikz}
\usepackage{pgfplots}

\newcommand{\wsnumber}{1}
\newcommand{\wstopic}{Vectors}
\pgfplotsset{
    every linear axis/.append style={
       axis x line=center,
       axis y line=center,
       xlabel={$x$},
       ylabel={$y$}
    },
    every axis plot/.append style={thick,mark=none}
}
\tikzset{
    point/.style={circle,draw,fill,minimum width=0.3ex,inner sep=0pt,outer sep=0pt},
    every label/.append style={black}
}


\usepackage[margin=1in]{geometry}
\usepackage{amsmath, amssymb, amsthm, graphicx, hyperref}
\usepackage{enumerate}
\usepackage{fancyhdr}
\usepackage{multirow, multicol}
\usepackage{tikz}
\pagestyle{fancy}
\fancyhead[RO]{Dennis Li}
\fancyhead[LO]{MA-UY 3113 Complex Variables and Linear Algebra }
\usepackage{comment}
\newif\ifshow
\showfalse

\ifshow
  \newenvironment{solution}{\textbf{Solution.}}{}
\else
  \excludecomment{solution}
\fi

\renewcommand{\thefootnote}{\fnsymbol{footnote}}
\usepackage{comment}


\newtheorem*{remark}{Remark}

\newcommand*{\GridSize}{4}

\newcommand*{\ColorCells}[1]{% #1 = list of x/y/color
  \foreach \x/\y/\color in {#1} {
    \node [fill=\color, draw=none, thick, minimum size=1cm] 
      at (\x-.5,\GridSize+0.5-\y) {};
    }%
}%

\begin{document}

\begin{center}
\ifshow
  \textbf{\Large Homework 8 Solution}\\
\else
  \textbf{\Large Homework 7}\\
\fi
Due: Monday April 29, by 11:59pm,\\via Gradescope\\
\end{center}

\hrule

\vspace{0.2cm}

\begin{enumerate}[$\bullet$]
\item  {\textbf{\textit{Note that you must assign a page to each problem you submit.}}}   Gradescope has great YouTube videos available on how to submit homework.  \textit{\textbf{Failure to submit homework correctly will result in a zero on homework.}}
\item No LaTeX for this homework.  
\item Late homework is not accepted.  Lateness due to technical issues will not be excused.  
\end{enumerate}

\hrule

\vspace{0.5cm}



\begin{enumerate}


 \item Orthogonally diagonalize the matrix $A$, whose eigenvalues are $7$ and $-2$. $A=\begin{bmatrix} 3 & -2 & 4 \\ -2 & 6 & 2 \\ 4 & 2 & 3 \end{bmatrix}$.  That is, find orthogonal matrix $Q$ and diagonal matrix $D$ such that $A = Q D Q^{T}$. 
\\
\textbf{Solution}\\
First, we would begin by finding the diagonalization of $A$, and we start from finding the eigenvectors
\[
\lambda = -2 \;\;\; 
A=\begin{bmatrix} 
5 & -2 & 4 \\ 
-2 & 8 & 2 \\ 
4 & 2 & 5 
\end{bmatrix}
\]
we reduce this matrix
\[
\begin{bmatrix}
5 & -2 & 4 \\
0 & \frac{36}{5} & \frac{18}{5} \\
0 & 0 & 0
\end{bmatrix}
\]
and we can obtain the eigenvector
\[
v_1 = \begin{bmatrix}
    -1\\
\end{bmatrix}
\]

 

\item Let $A$ be an $m \times n$ matrix.  Show that $AA^{T}$ and $A^{T}A$ have the same non-zero eigenvalues.  

\item A symmetric matrix is said to be positive semi-definite if and only if all of its eigenvalues are non-negative.  Show that 
\[
A  =  \begin{bmatrix} 1 & 2 & 1   \\ 2 & 4 & 2  \\ 1 & 2 & 1   \end{bmatrix} 
\]
is positive semi-definite.

\item Assume that $D$ is a diagonal matrix with non-negative diagonal entries.  We can define $\sqrt{D}$ to be the matrix that square roots the diagonal terms.  For instance, if 
\[
D = \begin{bmatrix} 4 & 0 & 0  \\ 0 & 9 & 0 \\ 0 & 0 & 0 \end{bmatrix} 
\]
then 
\[
\sqrt{D} = \begin{bmatrix} 2 & 0 & 0  \\ 0 & 3 & 0 \\ 0 & 0 & 0 \end{bmatrix} 
\]
If a symmetric matrix is positive semi-definite, then we can define $\sqrt{A} = Q \sqrt{D} Q^{T}$. Find $\sqrt{A}$ for 
\[
A  =  \begin{bmatrix} 1 & 2 & 1   \\ 2 & 4 & 2  \\ 1 & 2 & 1   \end{bmatrix} 
\]

\item Let 
\[
A = \begin{bmatrix}
    a & b \\
    b & c
\end{bmatrix}
\]
Assume that $|A| > 0$ and $a > 0$.  Show that $A$ is positive definite by showing that the eigenvalues of $A$ must be positive.  

\begin{remark}
Think back to a previous homework that relates the eigenvalues of a $2 \times 2$ matrix to its determinant and trace.  
\end{remark}
    


\item A function $f(x, y)$ has the following Taylor Expansion provided $(x, y) \sim (a, b)$
\[
\begin{aligned} 
f(x, y) \sim &f(a, b) + f_{x}(a, b) (x - a) + f_{y}(a, b) (y - b)  \\
&+ \frac{1}{2} f_{xx}(a,b) (x - a)^{2} + f_{xy}(a, b) (x - a)(y-b) + f_{yy}(a, b)(y-b)^{2} 
\end{aligned} 
\]
Assume $(a, b)$ is a critical value of $f$ such that $\nabla f(a, b) =0$.  Assume $|D(a, b)| > 0$ and $f_{xx}(a, b) > 0$ where 
\[
D(x, y) = \begin{bmatrix}
    f_{xx}(x, y) & f_{xy}(x, y) \\
    f_{xy}(x, y) & f_{yy}(x, y)
\end{bmatrix}
\]
Use the Taylor expansion given to show that $f(a, b)$ is a local minimum of $f$. 






\item Find the matrices $U, \;\Sigma, \; V$ in the SVD of $A=\begin{bmatrix} 1 & 5 & 2 \\ -3 & 1 & 2 \end{bmatrix}$.  

\begin{remark}
    When I ask for the SVD of a matrix whether on homework or an exam, I am referring to the reduced SVD.  That is, where the singular matrix $\Sigma$ has only positive singular values.  As a result, you will not need left and right singular vectors $\vec{u}, \; \vec{v}$ associated to the singular values equal to zero.  
\end{remark}

\item The image 
\[
\begin{tikzpicture}[scale=1]

    \begin{scope}[thick,local bounding box=name]
        \ColorCells{1/1/black, 2/3/black, 3/2/black, 1/4/black}
        \draw (0, 0) grid (\GridSize, \GridSize);
    \end{scope}

\end{tikzpicture}
\]
can be represented by a $4\times 4$ matrix $A$ of 1's and 0's by assigning 1 to any black square in the image and 0 to any white square in the image.  
\begin{enumerate}
    \item   Find the matrix $A$.  
    \item   The matrix $A$ will have 3 positive singular values, hence it will have an expansion of the form 
    \[
    A = \sigma_{1} \vec{u}_{1}\vec{v}_{1}^{T} + \sigma_{2} \vec{u}_{2}\vec{v}_{2}^{T}  + \sigma_{3} \vec{u}_{3}\vec{v}_{3}^{T} 
    \]
    Find the singular values and vectors in this expansion.  
    \item   Draw the image associated to the matrix $\sigma_{1} \vec{u}_{1} \vec{v}_{1}$, where $\sigma_{1}$ is the largest singular value.
\end{enumerate} 

we integrate over the PDF to obtain CDF as follows
\[
\int_{-\infty}^{\infty} f(x) \; dx
\]
first moment:
\[
\int_{-\infty}^{\infty} xf(x) \; dx
\]

        







\end{enumerate} 

\end{document}